% SOURCE docs/release/2026-10-29/source/docs/writing-repair/2026-10-29/src/25ac0d0d844f12b5.tex CHAR 0
\chapter{When a bank in the data ceases to be the same bank}
\label{gap29:funding}
A funding study follows institutions through time. Its economic interpretation
requires the deposit outcome and the exposure to refer to the same institution.
A fixed effect removes a persistent level difference; it cannot repair an outcome
assigned to the wrong institution. This distinction matters before a bank spends
money on a more elaborate covariance estimator.

\section{A controlled comparison}
The generated study follows 60 institutions for 20 dates. Exposure follows
$x_{i,t}=0.5x_{i,t-1}+v_{i,t}$, starting from zero, and the outcome is
\[
 y_{i,t}=\alpha_i+0.2\sin(t/3)+0.4x_{i,t}+u_{i,t}.
\]
The institution effects, exposure innovations and outcome disturbances are
independent standard normal draws. The target slope is therefore $0.4$ in this
specified economy. It is an illustrative response, not an estimated bank elasticity.
One alternative adds a common innovation with an institution-specific loading.
Another circularly reassigns the outcomes of five institutions after the tenth
date while leaving their exposures unchanged. The latter is an artificial
record error, not an authentic merger history.

There are 300 independent panels in each of the three cases. We compare a
pooled regression with a regression that removes institution and date means.
For the latter, the uncertainty calculations use observation-level HC2,
institution-clustered covariance, or the sum of institution and date cluster
covariances less the observation intersection. Reported intervals use the
normal critical value $1.96$. These are finite experiments with 20 dates, not
a general justification for that reference distribution.

\section{What changed the conclusion}
In clean panels, the pooled HC2 interval covers the target in $82.7\%$ of
requests. Within-institution and date adjustment gives $94.7\%$ with HC2 and
$94.3\%$ with institution clustering. Monte Carlo standard errors for these
coverage fractions are about 1.3--2.2 percentage points. Removing a nuisance
level and repairing a sampling calculation are different operations.

With reassigned outcomes, the within estimate has bias $-0.0170$, with Monte
Carlo standard error $0.00169$. Coverage is $90.7\%$ for HC2, $93.7\%$ for
institution clustering and $91.3\%$ for two-way clustering. The evaluator's
correct mapping gives bias $0.00102$ and institution-cluster coverage $95.0\%$.
The point is the remaining bias: a wider interval can sometimes cover a target
without restoring the meaning of the observation.

Under heterogeneous common shocks, the three within-regression coverage
fractions are $96.3\%$, $96.0\%$ and $94.0\%$. This construction does not show
that two-way clustering is universally preferable. One two-way variance is
nonpositive and is reported unavailable; it is included as a noncoverage in
the denominator of 300 requests. It is never clipped to manufacture an interval.

\section{The information worth obtaining}
An institution-history audit has a different value from another standard-error
option. A supplied warning identifies all 300 deliberately corrupted panels
and prompts refusal; this says nothing about a system's ability to discover an
unknown break. The next useful evidence would join exposures, outcomes and
institution histories using externally checked identifiers, then test whether
the estimated response survives the corrected economic unit.

The calculation was checked against an explicit regression with institution
and date indicator variables. Coefficients and all three covariance formulas
agree to better than $10^{-10}$ on the independent small-panel check. This
supports the reported arithmetic. It does not validate an unobserved bank's
history, guarantee finite-panel coverage, or establish a preferred estimator
outside the specified cases.

